\markdownRendererDocumentBegin
\markdownRendererSectionBegin
\markdownRendererHeadingOne{Hostile decoder substitution}\markdownRendererInterblockSeparator
{}The central alternative explanation is that the universal representation is penalized only because the downstream decoder has the wrong inductive bias. If so, stronger decoder-side search should buy back the compilation advantage. P11 treats that prediction as a mechanism test.\markdownRendererInterblockSeparator
{}\markdownRendererSectionBegin
\markdownRendererHeadingTwo{P11D sparse decoder — permanent negative}\markdownRendererInterblockSeparator
{}P11D preregistered a strong hostile gate: an L1 sparse universal decoder should still leave at least a 4× sample-threshold advantage for compiled state in both high-dimensional cells. It did not.\markdownRendererInterblockSeparator
{}\markdownRendererTable{}{3}{5}{drrrr}{{cell \markdownRendererCodeSpan{(d,s,r)}}{sparse universal \markdownRendererCodeSpan{n} at 0.95}{compiled \markdownRendererCodeSpan{n} at 0.95}{ratio}{compiled - sparse at \markdownRendererCodeSpan{n=64}}}{{(17,4,5)}{128}{64}{2×}{+0.2903}}{{(19,3,7)}{256}{64}{4×}{+0.3840}}\markdownRendererInterblockSeparator
{}The terminal \markdownRendererCodeSpan{P11D\markdownRendererUnderscore{}SPARSE\markdownRendererUnderscore{}DECODER\markdownRendererUnderscore{}GAP\markdownRendererUnderscore{}NOT\markdownRendererUnderscore{}MET} remains permanently negative. P11D also exposed an unseeded \markdownRendererCodeSpan{liblinear} replay defect; that defect is retained.\markdownRendererInterblockSeparator
{}
\markdownRendererSectionEnd \markdownRendererSectionBegin
\markdownRendererHeadingTwo{P11E deterministic sparse replication}\markdownRendererInterblockSeparator
{}P11E uses a fresh data seed and explicit estimator seeds. It reproduces thresholds 128/64 and 256/64 with low-sample compiled-minus-sparse advantages +0.2912 and +0.3307. Two fresh executions produce canonical SHA-256 \markdownRendererCodeSpan{1097d94bef1132d4dfa5d01176a9fcfcfebc46de8113e7cb2e57da1e579a4536}.\markdownRendererInterblockSeparator
{}
\markdownRendererSectionEnd \markdownRendererSectionBegin
\markdownRendererHeadingTwo{P11C/P11F/P11G nonlinear sequence}\markdownRendererInterblockSeparator
{}P11C's first execution attempt exceeded the available window; after an amendment that vectorized only its parity-bank evaluation, \markdownRendererCodeSpan{P11C\markdownRendererUnderscore{}EXECUTION\markdownRendererUnderscore{}RECEIPT\markdownRendererUnderscore{}V1.md} records the frozen protocol run to completion twice at \markdownRendererCodeSpan{P11C\markdownRendererUnderscore{}STRONGER\markdownRendererUnderscore{}DECODER\markdownRendererUnderscore{}GAP\markdownRendererUnderscore{}SUPPORTED}, with its pooled ≥4× gate passing at exactly the boundary and a twenty-seed sweep putting that boundary at 11 of 20 draws. It settles nothing and carries no claim authority; what it does establish is that its pooled combination rule was applied inside its own protocol. P11F produced positive numerical separation but is non-authoritative because hostile review found \markdownRendererCodeSpan{n\markdownRendererUnderscore{}jobs=-1} violated the frozen otherwise-default tree contract. P11G was frozen separately with \markdownRendererCodeSpan{n\markdownRendererUnderscore{}jobs=1}, explicit random states, and two fresh subprocess replays in the terminal decision path.\markdownRendererInterblockSeparator
{}\markdownRendererTable{}{3}{5}{drrrr}{{cell}{deterministic tree universal \markdownRendererCodeSpan{n} at 0.95}{compiled \markdownRendererCodeSpan{n} at 0.95}{tree accuracy at \markdownRendererCodeSpan{n=1024}}{compiled - tree at \markdownRendererCodeSpan{n=64}}}{{(17,4,5)}{\markdownRendererCodeSpan{NOT\markdownRendererUnderscore{}REACHED}}{64}{0.8248}{+0.4624}}{{(19,3,7)}{\markdownRendererCodeSpan{NOT\markdownRendererUnderscore{}REACHED}}{64}{0.7828}{+0.3942}}\markdownRendererInterblockSeparator
{}P11G's terminal is \markdownRendererCodeSpan{P11G\markdownRendererUnderscore{}DETERMINISTIC\markdownRendererUnderscore{}TREE\markdownRendererUnderscore{}DECODER\markdownRendererUnderscore{}GAP\markdownRendererUnderscore{}SUPPORTED}; both fresh scientific payloads have SHA-256 \markdownRendererCodeSpan{a2b0c33ce3c39e54ca1aa400a2b7d52d019fc4503f6cd5eb726c7b8bbe79a7cc}.\markdownRendererInterblockSeparator
{}
\markdownRendererSectionEnd \markdownRendererSectionBegin
\markdownRendererHeadingTwo{What P11G's terminal is a statement about}\markdownRendererInterblockSeparator
{}The programme registered three universal-state arms in P11C — \markdownRendererCodeSpan{UNIVERSAL\markdownRendererUnderscore{}L2}, \markdownRendererCodeSpan{UNIVERSAL\markdownRendererUnderscore{}L1} and \markdownRendererCodeSpan{UNIVERSAL\markdownRendererUnderscore{}EXTRA\markdownRendererUnderscore{}TREES} — and P11G's receipt publishes curves for one. Replaying P11G's own frozen data stream with only the decoder swapped, and reading P11G's own four scientific gates on each arm:\markdownRendererInterblockSeparator
{}\markdownRendererTable{}{4}{3}{ddd}{{universal arm}{0.95 threshold per cell, censored at 256}{terminal P11G's own gates print}}{{\markdownRendererCodeSpan{UNIVERSAL\markdownRendererUnderscore{}L2}}{≥256, ≥256}{\markdownRendererCodeSpan{...\markdownRendererUnderscore{}GAP\markdownRendererUnderscore{}SUPPORTED}}}{{\markdownRendererCodeSpan{UNIVERSAL\markdownRendererUnderscore{}L1}}{\markdownRendererStrongEmphasis{128}, ≥256}{\markdownRendererCodeSpan{...\markdownRendererUnderscore{}GAP\markdownRendererUnderscore{}NOT\markdownRendererUnderscore{}MET}}}{{\markdownRendererCodeSpan{UNIVERSAL\markdownRendererUnderscore{}EXTRA\markdownRendererUnderscore{}TREES} (reported)}{≥256, ≥256}{\markdownRendererCodeSpan{...\markdownRendererUnderscore{}GAP\markdownRendererUnderscore{}SUPPORTED}}}\markdownRendererInterblockSeparator
{}Two of three comparable pairs change the verdict, so the arm axis is not inert, and the flip is entirely the threshold gate: 128 is not ≥256. This is the same sparse threshold P11D reports as a permanent negative and P11E replicates; what is new is that placing it in P11G's gate, on P11G's own bytes, prints \markdownRendererCodeSpan{NOT\markdownRendererUnderscore{}MET}. \markdownRendererCodeSpan{NOT\markdownRendererUnderscore{}REACHED} through \markdownRendererCodeSpan{n=1024} is therefore not a stronger reading than 128 — an arm that reaches nothing anywhere gives the same gate reading in every world.\markdownRendererInterblockSeparator
{}P11G's terminal is retained exactly as frozen and is evidence about the decoder its own claim-authority sentence names. \markdownRendererCodeSpan{P11G\markdownRendererUnderscore{}ARM\markdownRendererUnderscore{}PLACEMENT\markdownRendererUnderscore{}ADJUDICATION\markdownRendererUnderscore{}V1.md} carries the adjudication, including the finding — read off the two freezes — that P11C's pooled combination rule governs P11C and does not bind P11G.\markdownRendererInterblockSeparator
{}
\markdownRendererSectionEnd \markdownRendererSectionBegin
\markdownRendererHeadingTwo{Decomposing a gap that moves two things at once}\markdownRendererInterblockSeparator
{}P11G compares L2 logistic regression on \markdownRendererCodeSpan{r} compiled columns against a 96-tree ensemble on the complete bank. Holding the decoder at ExtraTrees and moving only the representation separates them.\markdownRendererInterblockSeparator
{}\markdownRendererTable{}{3}{5}{drrrr}{{cell}{published gap at \markdownRendererCodeSpan{n=64}}{decoder-family half}{state half}{state share}}{{(17,4,5)}{+0.4624}{+0.0614}{+0.4010}{86.7\markdownRendererPercentSign{}}}{{(19,3,7)}{+0.3942}{+0.1757}{+0.2185}{55.4\markdownRendererPercentSign{}}}\markdownRendererInterblockSeparator
{}It cuts both ways and is reported both ways. It narrows the terminal: 13.3\markdownRendererPercentSign{} and 44.6\markdownRendererPercentSign{} of the published gaps are the change of decoder family, not of state. It supports the placement claim: with the decoder held fixed the state half is the majority in both cells, and being measured at a fixed decoder it is unaffected by which universal arm sits in the gate.\markdownRendererInterblockSeparator
{}The sequence supports the interpretation that \markdownRendererStrongEmphasis{compilation and decoder inductive bias are alternative locations for structural search}. Stronger downstream structure discovery should shrink the upstream advantage; the sparse negative is therefore part of the mechanism evidence, not an inconvenient result to erase. Each verdict in it is scoped to the arm that produced it.\markdownRendererInterblockSeparator
{}
\markdownRendererSectionEnd \markdownRendererSectionBegin
\markdownRendererHeadingTwo{P11H pooled successor — the survival was not predetermined, and it did not hold}\markdownRendererInterblockSeparator
{}The arm-scoping above is the smaller finding. The larger one, from \markdownRendererCodeSpan{orion.study.p11.attack\markdownRendererUnderscore{}audit}, is that all four of P11G's scientific gates hold in \markdownRendererStrongEmphasis{every world its own freeze admits}, so its survival was fixed before its seed was drawn. P11H re-asks the question under a protocol whose attack can win, editing nothing of P11G.\markdownRendererInterblockSeparator
{}P11H registers all three universal arms and freezes the best-of-arms rule inside its own positive gate; carries P11G's \markdownRendererCodeSpan{0.95} and \markdownRendererCodeSpan{0.20} thresholds over unedited; and draws two protected regimes by its fresh seed from a frozen 2×3 ladder of state widths \markdownRendererCodeSpan{r ∈ \markdownRendererLeftBrace{}3, 7\markdownRendererRightBrace{}} crossed with complete parity banks of 91, 364 and 969 columns. Before execution, the recorded preflight reports both hypothesis gates \markdownRendererCodeSpan{BOTH\markdownRendererUnderscore{}OUTCOMES\markdownRendererUnderscore{}REACHABLE} — supports \markdownRendererCodeSpan{[0.8808, 1.0000]} against \markdownRendererCodeSpan{0.95} and \markdownRendererCodeSpan{[0.0000, 0.2482]} against \markdownRendererCodeSpan{0.20} — and two reachable terminals over 15 admissible draws, 3 of which clear every gate.\markdownRendererInterblockSeparator
{}The seed drew \markdownRendererCodeSpan{(14,2,3)} and \markdownRendererCodeSpan{(19,3,3)}. Every precondition held, replay was byte-identical, and both hypothesis gates failed: \markdownRendererCodeSpan{P11H\markdownRendererUnderscore{}POOLED\markdownRendererUnderscore{}UNIVERSAL\markdownRendererUnderscore{}ATTACK\markdownRendererUnderscore{}PREVAILED}.\markdownRendererInterblockSeparator
{}\markdownRendererTable{}{7}{4}{drrr}{{rung}{universal bank}{pooled 0.95 threshold}{\markdownRendererCodeSpan{delta64}}}{{(14,2,3)}{91}{\markdownRendererStrongEmphasis{128}}{+0.1482}}{{(14,3,3)}{364}{\markdownRendererStrongEmphasis{128}}{+0.0992}}{{(19,3,3)}{969}{\markdownRendererStrongEmphasis{128}}{+0.0506}}{{(14,2,7)}{91}{≥256}{+0.2350}}{{(14,3,7)}{364}{≥256}{+0.3172}}{{(19,3,7)}{969}{≥256}{+0.3175}}\markdownRendererInterblockSeparator
{}The pool reaches the target by \markdownRendererCodeSpan{n=128} at every \markdownRendererCodeSpan{r=3} rung and at no \markdownRendererCodeSpan{r=7} rung, while the complete bank moves 91 → 969 columns inside each half without changing a verdict. The advantage is governed by the width of the compiled state, not by the size of the universal representation.\markdownRendererInterblockSeparator
{}And the decomposition sharpens in the opposite direction: at the drawn regimes the decoder-family half of the \markdownRendererCodeSpan{n=64} gap is exactly \markdownRendererCodeSpan{+0.0000}, so \markdownRendererStrongEmphasis{100\markdownRendererPercentSign{}} of it is the change of state — against 86.7\markdownRendererPercentSign{} and 55.4\markdownRendererPercentSign{} above — and the gap is still only \markdownRendererCodeSpan{+0.0506}. Attribution and magnitude are different questions.\markdownRendererInterblockSeparator
{}Three of the fifteen admissible draws would have printed the positive terminal and the seed did not draw one, so the \markdownRendererCodeSpan{r=7} rungs carry no terminal and no claim authority. The \markdownRendererCodeSpan{r=5} boundary was excluded before execution for verdict instability across three preflight seeds, in both directions, and remains open.\markdownRendererInterblockSeparator
{}
\markdownRendererSectionEnd \markdownRendererSectionBegin
\markdownRendererHeadingTwo{P11I wide high-width replication}\markdownRendererInterblockSeparator
{}P11I prospectively freezes the narrower regime P11H located rather than relabeling P11H. It evaluates the complete cross of three fresh execution seeds and three fixed bank-geometry strata at \markdownRendererCodeSpan{r=7}, with a matched \markdownRendererCodeSpan{r=3} control in every cell. The independent random unit is the execution seed (\markdownRendererCodeSpan{n=3}); geometry is a fixed within-seed stratum and five query repeats are technical repeats within each cell. One failed cell still defeats the conjunction.\markdownRendererInterblockSeparator
{}All nine high-width units pass. Compiled accuracy at \markdownRendererCodeSpan{n=64} ranges 0.9690–0.9981; the pooled attack's best accuracy below \markdownRendererCodeSpan{n=256} ranges 0.8489–0.9421 against the strict \markdownRendererCodeSpan{<0.95} gate; and \markdownRendererCodeSpan{delta64} ranges +0.2463–+0.3543 against \markdownRendererCodeSpan{>=+0.20}. The same pooled attack reaches 1.0000 in all nine matched \markdownRendererCodeSpan{r=3} controls. Two fresh subprocess payloads are byte-identical at SHA-256 \markdownRendererCodeSpan{b50ace30…e0ce}.\markdownRendererInterblockSeparator
{}The P11I terminal is \markdownRendererCodeSpan{P11I\markdownRendererUnderscore{}HIGH\markdownRendererUnderscore{}WIDTH\markdownRendererUnderscore{}ADVANTAGE\markdownRendererUnderscore{}REPLICATED\markdownRendererUnderscore{}WIDE\markdownRendererUnderscore{}PANEL}. It licenses a width-conditioned result only: the pooled attack wins in the narrow regime and the compiled-state advantage replicates in the registered high-width regime. A fresh two-subprocess revalidation under \markdownRendererCodeSpan{P11I\markdownRendererUnderscore{}REPLICATION\markdownRendererUnderscore{}UNIT\markdownRendererUnderscore{}AMENDMENT\markdownRendererUnderscore{}V1\markdownRendererUnderscore{}1.md} reproduces every cell byte-identically while recording \markdownRendererCodeSpan{n=3}, not nine. P11D and P11H remain adverse historical results.
\markdownRendererSectionEnd 
\markdownRendererSectionEnd \markdownRendererDocumentEnd